\section{Partie théorique : Ondes et interférences}
\label{sec:theorie}

Une onde indéformable est décrite par l'équation suivante, où \(x\) représente la composante spatiale et \(t\) la composante temporelle :

\begin{equation}
    p(x,t) = A \cdot \sin(\omega t - kx)
    \label{eq:onde}
\end{equation}

Cette onde est caractérisée par le nombre d'onde \(k = \frac{2\pi}{\lambda}\) [rad/m] et la pulsation \(\omega = 2\pi f\) [rad/s], où \(\lambda\) est la longueur d'onde [m] et \(f\) la fréquence [Hz]. L'amplitude \(A\) définit l'intensité maximale de l'onde.
\\\\
Cette section aborde deux cas spécifiques : les interférences d'ondes acoustiques dans un milieu non dispersif (l'air) et les interférences d'ondes de surface circulaires dans un milieu potentiellement dispersif (l'eau).

\subsection{Interférences d'ondes indéformables}
\label{subsec:interferences}

La superposition de deux ondes cohérentes \(p_1(x,t)\) et \(p_2(x,t)\) de même amplitude génère une onde stationnaire décrite par :

\begin{equation}
    p_m(t) = 2A \cos\left(k \frac{d_2 - d_1}{2}\right) \sin\left(\omega t - k \frac{d_1 + d_2}{2}\right)
    \label{eq:stationnaire}
\end{equation}

Les interférences sont constructives lorsque la différence de chemin \(\Delta x = d_2 - d_1\) est un multiple entier de la longueur d'onde \(\lambda\) :

\begin{equation}
    \Delta x = n \lambda \quad \text{avec } n \in \mathbb{Z}
    \label{eq:constructive}
\end{equation}

Elles sont destructives lorsque cette différence est un multiple demi-entier de \(\lambda\) :

\begin{equation}
    \Delta x = \left(n + \frac{1}{2}\right) \lambda \quad \text{avec } n \in \mathbb{Z}
    \label{eq:destructive}
\end{equation}

Ces principes s'appliquent aussi bien aux ondes acoustiques qu'aux ondes de surface bien que les motifs géométriques diffèrent selon la dimensionnalité du milieu.

\subsection{Ondes acoustiques}
\label{subsec:acoustique}

Les ondes acoustiques dans l'air suivent l'équation générale d'une onde indéformable. La célérité \(c\) de ces ondes est donnée par le rapport entre la pulsation \(\omega\) et le nombre d'onde \(k\) :

\begin{equation}
    c = \frac{\omega}{k}
    \label{eq:celerite}
\end{equation}

Cette relation peut également s'exprimer en fonction de la longueur d'onde \(\lambda\) et de la période \(T\) :

\begin{equation}
    c = \frac{\lambda}{T}
    \label{eq:celerite2}
\end{equation}

Dans l'air à pression atmosphérique, la célérité peut être calculée précisément par la formule :

\begin{equation}
    c = \sqrt{\frac{\gamma R T}{M}}
    \label{eq:celerite_exacte}
\end{equation}

où \(\gamma\) est le coefficient adiabatique, \(R\) la constante universelle des gaz parfaits, \(T\) la température absolue et \(M\) la masse molaire de l'air. Cette formule peut être approximée pour des températures proches des conditions standards (\(\theta\) en °C) par :

\begin{equation}
    c \approx 331 + 0.6 \cdot \theta \quad \text{[m/s]}
    \label{eq:celerite_approx}
\end{equation}

Cette approximation est valable pour un milieu non dispersif où la célérité ne dépend pas de la fréquence, mais seulement dans l'air à pression atmosphèrique.

\subsection{Ondes de surface circulaires}
\label{subsec:surface}

Les ondes de surface circulaires, contrairement aux ondes acoustiques, se propagent dans un milieu potentiellement dispersif comme l'eau. Leur équation en régime harmonique s'écrit tel que :

\begin{equation}
    h(r,t) = A(r) \cdot \sin(\omega t - kr)
    \label{eq:onde_surface}
\end{equation}

où \(A(r)\) est une fonction de Bessel dépendant de la distance radiale \(r\). La célérité \(c\) de ces ondes dépend de la profondeur \(d\) et de la longueur d'onde \(\lambda\). Cela implique donc que nous nous situons en un milieu dispersif :

\begin{equation}
    c(\lambda) = \sqrt{\left(\frac{g \lambda}{2\pi} + \frac{2\pi \gamma}{\rho \lambda}\right) \tanh\left(\frac{2\pi d}{\lambda}\right)}
    \label{eq:celerite_surface}
\end{equation}

avec \(g\) l'accélération gravitationnelle, \(\gamma\) la tension superficielle du liquide, et \(\rho\) sa masse volumique. En eau peu profonde (\(d \ll \lambda\)), la célérité peut être donnée par \(c = \sqrt{g d}\) par approximation quand d est beaucoup plus petit que la longueur d'onde.
\\\\
En isolant $\gamma$, nous obtenons :
\begin{equation}
    \gamma = \frac{\rho \lambda}{2\pi} \left( \frac{c^2}{\tanh\left(\frac{2\pi d}{\lambda}\right)} - \frac{g \lambda}{2\pi} \right)
    \label{eq:gamma_isole}
\end{equation}

Les interférences entre deux sources d'ondes de surface forment alors des motifs hyperboliques qui reflétant en image (voir Fig. \ref{fig:intro}) la nature bidimensionnelle de leur propagation.
